% Réplica con cambio manual (ADR-0002): funcionamiento normal, fallo de la
% computadora y regreso de la computadora como réplica.
\documentclass[tikz,border=0pt]{standalone}
% Estilo común de los diagramas de los ADR.
\usepackage[T1]{fontenc}
\usepackage[utf8]{inputenc}
\usepackage[sfdefault]{FiraSans}
\usepackage{fontawesome5}
\usetikzlibrary{positioning,backgrounds,fit,arrows.meta,calc}

% Capa para los recuadros de zona: encima del fondo blanco y debajo de los equipos.
\pgfdeclarelayer{zonas}
\pgfsetlayers{background,zonas,main}

\definecolor{texto}{HTML}{262626}
\definecolor{azul}{HTML}{1F4E79}
\definecolor{naranja}{HTML}{C55A11}
\definecolor{gris}{HTML}{8C8C8C}

\tikzset{
  % Se aplica en cada tikzpicture: fondo blanco para que se lea en modo oscuro.
  diagrama/.style={
    font=\small, color=texto,
    background rectangle/.style={fill=white},
    show background rectangle, inner frame sep=14pt,
  },
  % Dispositivo con icono, título y descripción opcional.
  equipo/.style={
    draw=#1, fill=#1!10, line width=0.9pt, rounded corners=5pt,
    align=center, font=\footnotesize, inner sep=6pt, minimum width=3.1cm, minimum height=1.9cm,
  },
  equipo/.default=azul,
  % Recuadro que agrupa los equipos de un centro o una etapa.
  zona/.style={
    draw=#1, line width=0.8pt, rounded corners=8pt, inner sep=12pt,
  },
  zona/.default=gris,
  titulozona/.style={
    anchor=south west, font=\small\bfseries, text=#1, inner sep=3pt,
  },
  titulozona/.default=gris,
  enlace/.style={Latex-Latex, line width=0.9pt, draw=azul},
  copia/.style={-Latex, line width=1.1pt, draw=naranja},
  etiqueta/.style={font=\footnotesize, fill=white, inner sep=2pt, align=center},
  nota/.style={
    draw=#1, fill=#1!6, line width=0.7pt, rounded corners=4pt,
    font=\footnotesize, align=center, inner sep=5pt,
  },
  nota/.default=gris,
}

% \equipotexto{icono}{título}{descripción}
% La descripción usa el tamaño de letra del nodo, porque TikZ no conserva un cambio
% de tamaño después de \\ ni admite \\ dentro de un grupo.
\newcommand{\equipotexto}[3]{{\Large #1}\\[3pt]{\small\bfseries #2}\\[1pt]#3}


% \telefonos{prefijo}{nodo servidor}: tres voluntarios conectados al servidor.
\newcommand{\telefonos}[2]{
  \foreach \i/\x in {1/0.3, 2/2.15, 3/4.0} {
    \node[equipo, minimum width=1.7cm, minimum height=1.3cm] (#1t\i) at (\x,-3.3)
      {{\Large\faMobile*}\\[2pt]Voluntario};
    \draw[enlace] (#2.south) -- (#1t\i.north);
  }
}

\begin{document}
\begin{tikzpicture}[diagrama]
  \tikzset{equipo/.append style={minimum width=2.9cm}}

  % 1. Funcionamiento normal.
  \begin{scope}
    \coordinate (arriba1) at (2.15,1.55);
    \node[equipo] (a1) at (0,0)
      {\equipotexto{\faLaptop}{Servidor principal}{Computadora\\Turno 1}};
    \node[equipo=naranja] (b1) at (4.3,0)
      {\equipotexto{\faIcon{tablet-alt}}{Réplica}{Tablet del coordinador\\Turno 1}};
    \draw[copia] (a1.east) -- (b1.west);
    \node[etiqueta, text=naranja, anchor=south] at ($(a1.north east)!0.5!(b1.north west)$) {Copia cada cambio};
    \telefonos{p1}{a1}
    \node[nota] (n1) at (2.15,-5.1) {Los voluntarios trabajan con la computadora.\\Cada registro se copia en la tablet.};
    \begin{pgfonlayer}{zonas}
      \node[zona, fit=(arriba1) (a1) (b1) (p1t1) (p1t3) (n1)] (z1) {};
    \end{pgfonlayer}
    \node[titulozona=texto] at (z1.north west) {1. Funcionamiento normal};
  \end{scope}

  % 2. Se apaga la computadora.
  \begin{scope}[xshift=8.7cm]
    \coordinate (arriba2) at (2.15,1.55);
    \node[equipo=gris, dashed, text=gris] (a2) at (0,0)
      {\equipotexto{\faPowerOff}{Computadora}{Apagada\\o sin batería}};
    \node[equipo] (b2) at (4.3,0)
      {\equipotexto{\faIcon{tablet-alt}}{Nuevo servidor}{Tablet del coordinador\\Turno 2}};
    \telefonos{p2}{b2}
    \node[nota=azul] (n2) at (2.15,-5.1) {\faHandPointer\ El coordinador pulsa «Convertir en servidor».\\Los voluntarios pasan a trabajar con la tablet.};
    \begin{pgfonlayer}{zonas}
      \node[zona, fit=(arriba2) (a2) (b2) (p2t1) (p2t3) (n2)] (z2) {};
    \end{pgfonlayer}
    \node[titulozona=texto] at (z2.north west) {2. Se apaga la computadora};
  \end{scope}

  % 3. Vuelve la computadora.
  \begin{scope}[xshift=17.4cm]
    \coordinate (arriba3) at (2.15,1.55);
    \node[equipo=naranja] (a3) at (0,0)
      {\equipotexto{\faLaptop}{Réplica}{Computadora\\Turno 2}};
    \node[equipo] (b3) at (4.3,0)
      {\equipotexto{\faIcon{tablet-alt}}{Servidor}{Tablet del coordinador\\Turno 2}};
    \draw[copia] (b3.west) -- (a3.east);
    \node[etiqueta, text=naranja, anchor=south] at ($(a3.north east)!0.5!(b3.north west)$) {Se pone al día};
    \telefonos{p3}{b3}
    \node[nota] (n3) at (2.15,-5.1) {La computadora ve un turno mayor que el suyo,\\deja de aceptar registros y pasa a ser réplica.};
    \begin{pgfonlayer}{zonas}
      \node[zona, fit=(arriba3) (a3) (b3) (p3t1) (p3t3) (n3)] (z3) {};
    \end{pgfonlayer}
    \node[titulozona=texto] at (z3.north west) {3. Vuelve la computadora};
  \end{scope}

\end{tikzpicture}
\end{document}
